\chapter{Conclusion}

This thesis studied Legal Judgment Prediction for Indian court cases as a
structured prediction problem rather than a plain document-classification
task. It argued that credible legal prediction requires more than strong
accuracy numbers: it requires a representation that preserves legally
meaningful structure while actively resisting leakage from decision-bearing
text and identity-based shortcuts.

To address this, the thesis introduced an end-to-end pipeline that converts
raw Indian judgments into leakage-audited, graph-ready case records through
entity extraction, rhetorical-role segmentation, outcome labelling,
multi-hearing consolidation, and reasoning-focused graph construction. On
top of this representation, a Heterogeneous Graph Transformer was used to
model both the internal structure of each case and the shared legal links
across cases through statutes, provisions, and precedents.

Across more than 73{,}000 Indian court cases spanning five legal domains,
the proposed framework achieved strong pooled performance
($0.8020$ accuracy, $0.7959$ macro-F1, $0.8831$ ROC-AUC under 5-fold
cross-validation on the best \texttt{party\_args\_lr\_decay}
configuration) and consistently outperformed text-only reference settings.
More importantly, the results and ablations suggest that the gains are not
coming merely from surface-level text patterns, but from cross-case legal
structure encoded in a controlled and auditable way: a probing classifier
isolates a $+3.97$-point accuracy gain attributable to graph message passing
at fixed classifier capacity.

Beyond pooled accuracy, the thesis addressed the two questions a legal
audience actually cares about: \emph{can individual predictions be
read}, and \emph{does the model transfer beyond the distribution it was
trained on}? A four-level interpretability pipeline --- representation
probing, component-level importance, error decomposition, and the
case-level \path{Graph_Analyser} built around a heterogeneous
PGExplainer, FAISS retrieval in post-GNN space, and LLM verbalisation ---
produces grounded, legally named explanations for individual predictions
and exposes a clean structural signature of the model's confident
failures: when the prediction is confidently wrong, the retrieved
representation-space neighbours almost always reinforce the \emph{wrong}
label. A zero-shot evaluation on $11{,}769$ food-safety judgments, a
legal domain absent from training, reaches $0.608$ accuracy and
$0.675$ ROC-AUC against a $0.543$ majority baseline without any
fine-tuning, showing that the cross-bucket HGT learns a partially
domain-general legal-reasoning prior rather than memorising training
domains' statutory hubs.

The main contribution of this thesis is therefore methodological. It shows
that legal outcome prediction over Indian judgments becomes more credible
when the task is built around leakage-safe data construction, heterogeneous
graph design, explicit constraints on what information may be shared
across cases, per-case explainability tied to the same representation the
model predicts from, and an out-of-domain evaluation that tests whether
the learned signal is truly about legal reasoning or about corpus
memorisation. While the improvements are measured rather than dramatic, they
support a clear conclusion: graph-based legal prediction is most useful when
it is not only accurate, but also structurally faithful, auditable at the
case level, and experimentally honest about its transfer behaviour.

Future work can build on this foundation by tightening leakage audits
further, extending the outcome space beyond binary decisions, designing
models that adapt to domain-specific difficulty and long factual
narratives, closing the in-domain/out-of-domain gap through light-weight
fine-tuning or domain-adaptive heads, and pairing the case-level
explainer with a legal-expert validation loop to convert it from an
internal diagnostic into an auditable review tool.